% ===== tr-06 块1：第 4 章 pp.57--60（原文页 57--60）=====
% [原文 p.57]
\figimg[0.72]{p4_fig09.png}
\figtitle{4.9}{$S(\mathfrak{D})$（$\ell=1$ 的情形）。这里各处的自同构都已明确给出，而恒等映射 $\theta_{\gamma,S(\mathfrak{D})}$ 正是这些给定自同构的复合。}
% [Tutor Note: 图 4.9 中标出 5 个自同构：直线 $\R(1,0)$ 与 $\R(0,1)$ 各被环道 $\gamma$ 穿过两次（左/右、上/下各得一个互逆的自同构），而新增射线 $\R_{\ge0}(1,1)$ 只被穿过一次，合计 $2+2+1=5$，正与正文（p.56 末）“这五个自同构的复合是恒等”相符。]

若 $\ell=2$，则人们发现
\begin{align*}
S(\mathfrak{D})\setminus\mathfrak{D}
&=\bigl\{(\R(n+1,n),(1+t^{2n+1}x^{-(n+1)}y^{-n})^{2})\bigm|n\in\Z,\ n\ge1\bigr\}\\
&\quad\cup\bigl\{(\R(n,n+1),(1+t^{2n+1}x^{-n}y^{-(n+1)})^{2})\bigm|n\in\Z,\ n\ge1\bigr\}\\
&\quad\cup\bigl\{(\R(1,1),(1-t^{2}x^{-1}y^{-1})^{-4})\bigr\}.
\end{align*}
这是我自己与 Siebert 最初通过计算机程序实验发现的，其第一个验证由 [14] 给出。它也立即由 [28] 的结果推出，这些结果将在下文中解释。
% [Tutor Note: 第三行末的指数是 $-4$（负四次幂），不是 $-1$；页图上 $-4$ 的“4”清晰可辨。前两行 $t$ 的指数 $2n+1$ 为奇数，与 $\ell=2$ 时附着函数取平方相对应。]

若 $\ell=3$，情形变得更加复杂。首先，正如 Kontsevich 注意到的，$S(\mathfrak{D})$ 具有一定的周期性。即
\[
(\R_{\ge0}(m_1,m_2),f(x^{-m_1}y^{-m_2}))\in S(\mathfrak{D})
\]
当且仅当
\[
(\R_{\ge0}(3m_1-m_2,m_1),f(x^{-(3m_1-m_2)}y^{-m_1}))\in S(\mathfrak{D}),
\]
只要 $m_1$、$m_2$ 与 $3m_1-m_2$ 都为正。此外，还有支撑为 $\R_{\ge0}(3,1)$ 与 $\R_{\ge0}(1,3)$ 的射线，于是由周期性，还有支撑为
\[
\R_{\ge0}(8,3),\ \R_{\ge0}(21,8),\ \dots\quad\text{与}\quad \R_{\ge0}(3,8),\ \R_{\ge0}(8,21),\ \dots
\]
的射线，它们收敛到斜率为 $(3\pm\sqrt5)/2$ 的射线，对应该线性变换
\[
\begin{pmatrix}3&-1\\1&0\end{pmatrix}
\]
的两个不同特征空间。这些射线中的每一条都具有形式
\[
(\R_{\ge0}(m_1,m_2),(1+t^{m_1+m_2}x^{-m_1}y^{-m_2})^{3}).
\]
% [Tutor Note: $(3\pm\sqrt5)/2$ 正是矩阵 $\begin{psmallmatrix}3&-1\\1&0\end{psmallmatrix}$ 的特征值：特征多项式 $\lambda^2-3\lambda+1=0$ 的根；而周期性变换 $(m_1,m_2)\mapsto(3m_1-m_2,m_1)$ 用的正是同一矩阵。页图上特征值的根式与矩阵元素均可逐字核对。]
% [原文 p.58]
这些是出现在由斜率 $(3\pm\sqrt5)/2$ 的两条射线所张成的锥之外的唯一一些射线。另一方面，在这个锥内部，每个有理斜率都出现，且所附的函数非常复杂。例如，附着于斜率 1 的直线上的函数是
\[
\left(\sum_{n=0}^{\infty}\frac{1}{3n+1}\binom{4n}{n}t^{2n}x^{-n}y^{-n}\right)^{9}.
\]
% [Tutor Note: 上式括号外的指数是 $9$，括号内系数为 $\frac{1}{3n+1}\binom{4n}{n}$；多项式系数 $\binom{4n}{n}$ 不是 $\binom{4n}{2n}$，页图上括号内上行为 $4n$、下行为 $n$。]

同样，这个形式是我与 Siebert 通过计算机实验发现的，但它由 Reineke 在 [69] 中作了验证。最近，Kontsevich 证明了附着于所有这些射线的函数都是代数的。例如，若 $g$ 表示上述函数的 9 次方根，则它满足方程
\[
t^{2}x^{-1}y^{-1}g^{4}-g+1=0.
\]

这一系列例子还与许多其他有趣的对象有联系。一方面，Reineke 在 [69] 中用 Kronecker $\ell$-箭图（即有两个顶点、两顶点之间有 $\ell$ 条箭头的箭图）的表示模空间的 Euler 示性数，对所附函数给出了解释。另一方面，这些图也与由这些箭图定义的 cluster 代数密切相关。这一联系将在与 Hacking、Keel 及 Kontsevich 的合作工作 [27] 中得到更详细的研究。

下面我们来说明 $S(\mathfrak{D})$ 中出现的函数的计数解释。为说明其动机，让我们回到第 4.4 节的热带解释。比如说，从一个对应于 K3 曲面的热带流形 $B$ 开始（如 图 4.10 所示），并配上我们将称之为\emph{热带圆盘}的东西。它几乎就是一条热带曲线，但它只是在点 $P$ 处终止，而在 $P$ 处没有任何平衡条件；同时，它的其他一些腿终止于 $B$ 的奇点。正如第 4.4 节末尾所解释的那样，这是合法的行为。按照第 4.4 节末尾的描述，可以设想在每个终止于奇点的腿上方有一个圆盘。在这些腿相交的地方，人们希望把这些圆盘粘合起来，并沿 $P$ 附近线段的柱面继续。终止于 $P$ 时，我们大致得到 $X(B)$ 中一个边界包含在 $P$ 上方环面纤维中的圆盘，如图所示。一个自然的问题是：初始的那些圆盘（可能对这些圆盘取多重覆盖）有多少种方式可以粘合起来，从而给出一个新的圆盘。

现在把这个图景与我们在镜像一侧所看到的情况作比较。我们的显式退化确实以 $\check X(B)$ 之类的东西作为一般纤维给出。然而，它由类似的热带信息所控制：射线从奇点沿单值不变的方向发出，正如热带曲线的情形一样。它们相撞，而命题 4.11.2 中 Kontsevich--Soibelman 的结果给出新的射线。因此人们可以期望，这一过程恰好反映了 $X(B)$ 中以 $X(B)\to B$ 的纤维为边界的全纯圆盘。
% [Tutor Note: $\check X(B)$ 的 check 帽在页图上清楚可辨（$X$ 上方有 $\check{}$ 记号）；这是 SYZ 文献中镜像流形的惯用记号，与 $X(B)$ 互为镜像，不要写成 $\bar X$ 或 $\tilde X$。]

还值得一提的是 Auroux 的工作 [4]，它更精确地刻画了如下观念：一侧的复结构应当由另一侧的全纯圆盘所决定。
% [原文 p.59]
\figimg[0.72]{p4_fig10.png}
\figtitle{4.10}{仿射 K3 曲面上的一个热带圆盘。这里 $\times$ 表示奇点，而该圆盘在点 $P$ 处“终止”。}
% [Tutor Note: 原文“the other.”是接 p.58 末句“... determined by holomorphic disks on”的收尾，本译本已按整句处理并译于上一段；图 4.10 位于该页顶部，先图后文。]

这也为如下想法提供了后验支持：在生成新射线的过程中必定存在一种计数解释。

处理全纯圆盘通常是困难的。把涉及全纯圆盘的问题转化为涉及真正的 Gromov--Witten 不变量的问题，往往更容易。对于这里所讨论的问题，我们可以这样做。于是，在 [28] 中所解释的最简单情形下，计数解释如下。

设给定互不相同的非零本原向量 $m_1,\dots,m_p\in M$ 以及正整数 $\ell_1,\dots,\ell_p$。考虑散射图
\[
\mathfrak{D}=\{(\R m_i,(1+tz^{-m_i})^{\ell_i})\mid 1\le i\le p\}.
\]
设 $(\mathfrak{d},f_\mathfrak{d})\in S(\mathfrak{D})\setminus\mathfrak{D}$。我们总可以假定它是 $S(\mathfrak{D})\setminus\mathfrak{D}$ 中唯一一条以 $\mathfrak{d}$ 为底空间集的射线。这是因为，若在 $S(\mathfrak{D})\setminus\mathfrak{D}$ 中有射线 $(\mathfrak{d}_1,f_{\mathfrak{d}_1}),(\mathfrak{d}_2,f_{\mathfrak{d}_2}),\dots$ 满足 $\mathfrak{d}_1=\mathfrak{d}_2=\cdots$，则我们可以把这一族射线替换为单条射线 $(\mathfrak{d}_1,\prod f_{\mathfrak{d}_i})$，而不影响 $\theta_{\gamma,S(\mathfrak{D})}$。在这一假定下，$f_\mathfrak{d}$ 由 $\mathfrak{D}$ 唯一确定。我们希望以计数方式解释 $f_\mathfrak{d}$。

为此，考虑 $M_\R$ 中一个完全扇 $\Sigma$，它的一维射线为
\[
\R_{\le0}m_1,\ \dots,\ \R_{\le0}m_p,\ \mathfrak{d}.
\]
为简化本讨论，设 $\mathfrak{d}$ 不与另外 $p$ 条射线重合。设 $X$ 是由 $\Sigma$ 定义的环面簇，其环面除子 $D_1,\dots,D_p,D_{\mathrm{out}}$ 对应于上述射线。接着，在除子 $D_i$ 上选取 $\ell_i$ 个一般位置的点，记为 $P_{i1},\dots,P_{i\ell_i}$。设 $\nu:\tilde X\to X$ 为对这些 $\sum_i\ell_i$ 个点的爆破，其例外除子 $E_{ij}$ 位于 $P_{ij}$ 之上。以 $\tilde D_i$、$\tilde D_{\mathrm{out}}$ 记 $D_i$、$D_{\mathrm{out}}$ 的真变换。
% [Tutor Note: 射线的方向与其函数中的单项式指数必须配套：$(\R m_i,(1+tz^{-m_i})^{\ell_i})$ 里射线方向是 $m_i$，而函数中出现的是 $z^{-m_i}$（负指数），页图上 $-m_i$ 的负号与下标 $i$ 均可核对。]

在下文中，我们将用记号 $\mathbf{P}_i=(p_{i1},\dots,p_{i\ell_i})$ 表示某个非负整数 $|\mathbf{P}_i|=p_{i1}+\cdots+p_{i\ell_i}$ 的一个长度为 $\ell_i$ 的分拆，其中允许某些 $p_{ij}$ 为零。固定
% [原文 p.60]
一个类 $\beta\in H^2(X,\Z)$，使得 $a_i:=\beta\cdot D_i$ 非负、$k:=\beta\cdot D_{\mathrm{out}}$ 为正。这是环面几何中的一个容易的练习：这蕴含着关系
\[
\sum_{i=1}^{p}a_im_i=km_{\mathrm{out}},
\]
其中 $m_{\mathrm{out}}$ 是 $\mathfrak{d}$ 的一个本原生成元。若选取一族分拆 $\mathbf{P}=(\mathbf{P}_1,\dots,\mathbf{P}_p)$，其中 $\mathbf{P}_i$ 是 $a_i$ 的一个分拆，则令
\[
\beta_{\mathbf{P}}:=\nu^*\beta-\sum_{i=1}^{p}\sum_{j=1}^{\ell_i}p_{ij}E_{ij}.
\]
这可以看作 $X$ 上一条曲线的类，该曲线恰好 $p_{ij}$ 次经过点 $P_{ij}$。
% [Tutor Note: 正文把 $\beta$ 取在 $H^2(X,\Z)$（上同调）中（页图可辨为上标 $2$），而把 $\beta_{\mathbf P}$ 称为“曲线的类”是靠 Poincaré 对偶实现的，不要误改成 $H_2$；另外分拆的记号是正体粗体 $\mathbf{P}$，与点 $P_{ij}$ 的斜体 $P$ 不同。]

我们想对这个上同调类配上一个数。这将是一个关于 $\tilde X$ 中带一个标记点的有理曲线的 Gromov--Witten 计数，这些曲线 (1) 代表类 $\beta_{\mathbf{P}}$；(2) 在标记点处以 $k$ 阶与 $\tilde D_{\mathrm{out}}$ 相切；(3) 除此之外与任一除子 $\tilde D_i$ 不相交。这是一个相对 Gromov--Witten 不变量。然而，经典的相对 Gromov--Witten 不变量理论是相对于光滑除子的，而这里的边界除子之并当然是奇异的。人们可以改用对数 Gromov--Witten 理论来编码上述条件。在撰写 [28] 时，对数 Gromov--Witten 理论尚不可用，因此我们使用了一个技术性的变通办法来化归到经典理论。我在这里给出这一描述，是因为它不需要了解对数 Gromov--Witten 理论。

人们定义 $\tilde X^{o}:=\tilde X\setminus\bigcup_{i=1}^{p}\tilde D_i$。然后考虑亏格零、目标空间为 $\tilde X^{o}$、相对于除子 $\tilde D^{o}_{\mathrm{out}}=\tilde D_{\mathrm{out}}\cap\tilde X^{o}$ 的相对稳定映射的模空间 $\mathfrak{M}(\tilde X^{o}/\tilde D^{o}_{\mathrm{out}},\beta_{\mathbf{P}})$。这些曲线有一个标记点，其与 $\tilde D^{o}_{\mathrm{out}}$ 的相切阶为 $k$。唯一的问题是目标空间非紧，但可以证明这不会造成任何问题，因为模空间仍然是紧的。人们发现它的虚拟维数为零，且由于它带有虚拟基本类，我们可以定义
\[
N_{\mathbf{P}}:=\int_{[\mathfrak{M}(\tilde X^{o}/\tilde D^{o}_{\mathrm{out}},\beta_{\mathbf{P}})]^{\mathrm{vir}}}1\in\Q.
\]
于是我们可以陈述计数结果（[28]）：

\begin{theorem}{4.11.4}
我们有
\[
\log f_\mathfrak{d}=\sum_{\beta}\sum_{\mathbf{P}}k(\beta)N_{\mathbf{P}}\,t^{\sum_i|\mathbf{P}_i|}z^{-k(\beta)m_{\mathrm{out}}},
\]
其中求和遍历所有满足 $\beta\cdot D_i\ge0$、$k(\beta):=\beta\cdot D_{\mathrm{out}}>0$ 的 $\beta\in H^2(X,\Z)$，以及满足 $|\mathbf{P}_i|=\beta\cdot D_i$ 的分拆 $\mathbf{P}$。
\end{theorem}
% [Tutor Note: 模空间的字体是 Fraktur 体的 $\mathfrak{M}$（页图可辨），不是花体的 $\mathcal{M}$；定理 4.11.4 中 $t$ 的指数写成 $\sum_i|\mathbf{P}_i|$（即各分拆长度之和），$z$ 的指数为 $-k(\beta)m_{\mathrm{out}}$，其中 $m_{\mathrm{out}}$ 是 $\mathfrak{d}$ 的本原生成元。]
